\documentclass[11pt,a4paper]{article}

\usepackage[margin=1in]{geometry}
\usepackage{amsmath,amssymb}
\usepackage{hyperref}
\usepackage{graphicx}
\usepackage{booktabs}
\usepackage{listings}
\usepackage{xcolor}
\usepackage{enumitem}

\hypersetup{colorlinks=true,linkcolor=blue,citecolor=blue,urlcolor=blue}

\lstset{
  basicstyle=\ttfamily\small,
  breaklines=true,
  frame=single,
  backgroundcolor=\color{gray!5},
}

\title{zkFetch: Fetch with a Proof}
\author{Anudit Nagar \\ \texttt{zkfetch}}
\date{October 2026}

\begin{document}
\maketitle

\begin{abstract}
\noindent
zkFetch is a protocol that turns an ordinary HTTPS request into verifiable web data.
A client fetches a URL together with a \emph{notary} and receives a normal HTTP response
plus a signed attestation.
Later, the client decides what to reveal --- a URL, headers, JSON fields --- or proves
facts about values it keeps hidden, such as \texttt{streak.length >= 365}.
Anyone can verify the resulting presentation offline, against a pinned notary key.
The notary never learns private data such as bearer tokens, and the client cannot
forge bytes the server did not send.
The protocol is built on TLSNotary with two session modes (MPC-TLS and proxy),
QuickSilver zero-knowledge proofs over authenticated TLS transcripts, an optional
offline Binius64 backend, and a \texttt{fetch}-shaped SDK.
This paper describes the full protocol as implemented: parties and trust model,
session flow, commitments, disclosure, verification, and measured performance.
\end{abstract}

\section{Introduction}

The web runs on TLS, but TLS authenticates a connection, not data.
Once a response is read, there is nothing to show a third party:
screenshots and copied JSON prove nothing.
Existing oracles solve this with trusted hardware, trusted servers, or by
revealing everything.

zkFetch takes a different approach.
It adds one party, the \emph{notary}, to a normal fetch:

\begin{center}
\texttt{zkFetch(url)} $\rightarrow$ (\texttt{Response}, attestation, secrets) \\
\texttt{present(spec)} $\rightarrow$ presentation \qquad
\texttt{verify(presentation)} $\rightarrow$ accept / reject
\end{center}

Three properties define the protocol:

\begin{enumerate}[leftmargin=2em]
\item \textbf{Server authenticity.} If verification succeeds, the disclosed bytes
and proven predicates describe what the named server actually sent.
\item \textbf{Privacy.} The notary learns ciphertext, record lengths, and
proof transcripts, but never plaintext, application keys, or hidden values.
The verifier learns only what the presentation discloses plus one bit per
predicate.
\item \textbf{Selectivity.} Disclosure is decided \emph{after} the fetch.
A session is stored, and each presentation reveals a subset: headers by name,
JSON fields by dotted path, the request target or body, or numeric facts
about hidden fields.
\end{enumerate}

No change to the server is required. No trusted hardware is used.
The prover runs on servers (native Rust), in browsers (single- and
multi-threaded WebAssembly), and in Chrome extensions; the notary is a small
stateless Rust service suitable for a single ARM machine.

\section{Parties, keys, and artifacts}

There are four roles. The \textbf{prover} $P$ is the client that fetches.
The \textbf{notary} $N$ co-signs what happened. The \textbf{server} $S$ is an
unmodified HTTPS origin. The \textbf{verifier} $V$ checks a presentation offline.

$N$ holds a long-term signing key (secp256k1 ECDSA by default); $P$ pins the
expected public key before connecting and aborts on mismatch.
$P$ and $N$ communicate over a WebSocket; $N$ rate-limits, authenticates
tenants by capability token, and enforces byte, record, session, and
predicate-work caps.

A session produces two artifacts:

\begin{itemize}[leftmargin=2em]
\item \textbf{Attestation.} $N$'s signature over: TLS version, cipher suite and
group; server name, dialed IP, port, certificate hashes, and whether $N$
validated the chain; transcript hashes; byte counts; hash commitments and
predicate results; session mode (\texttt{mpc} or \texttt{proxy}); wall-clock
time; and two binding labels, \texttt{owner} and \texttt{context}.
\item \textbf{Secrets.} Private openings held by $P$: plaintext, keys, and
commitment randomness. Secrets are treated as credentials; serialised sessions
must be stored encrypted.
\end{itemize}

A \textbf{presentation} combines the attestation with openings for exactly the
disclosed units plus zero-knowledge proofs for each claimed predicate.
Verification needs only the presentation, the pinned notary key, and a policy
(expected server, method, target, status, freshness window, challenge context,
required predicates).

The two binding labels deserve emphasis.
\texttt{context} carries a verifier-supplied nonce (challenge), which defeats
replay of old presentations.
\texttt{owner} carries an application-chosen label such as a wallet address;
it binds, but does not by itself prove ownership --- the application must still
check a separate wallet signature.

\section{Protocol overview}

At a high level every session has five phases:

\begin{enumerate}[leftmargin=2em]
\item \textbf{Open and admit.} $P$ opens a WebSocket to $N$, presents a
capability, checks $N$'s signed key-possession challenge against the pin, and
commits to mode, TLS version, and limits.
\item \textbf{TLS.} $P$ and $N$ jointly (MPC mode) or via $N$'s relay (proxy
mode) establish TLS with $S$ and exchange one HTTP request and response.
\item \textbf{Commit and prove.} $P$ commits to the transcript and proves
selected predicates in zero knowledge to $N$.
\item \textbf{Attest.} $N$ validates the proofs, then signs the attestation.
\item \textbf{Present and verify.} At any later time, $P$ builds a presentation
for a chosen disclosure spec; any $V$ verifies it offline.
\end{enumerate}

The SDK mirrors \texttt{fetch} deliberately:

\begin{lstlisting}[language=Java]
const res = await zkFetch("https://api.example.com/me", {
  headers: { Authorization: `Bearer ${token}` },
  zkConfig: {
    notaryUrl, expectedNotaryKey,
    mode: "proxy",
    tlsVersion: "auto",
    backend: "quicksilver",
    predicates: [{ jsonPath: "streak.length",
                   predicate: { gte: "365" } }],
    context: challenge,
  },
});
await res.json();
const presentation = res.zk.present({
  response: { jsonPaths: ["username"] },
  prove: [{ jsonPath: "streak.length",
            predicate: { gte: "365" } }],
});
const result = verify(presentation, {
  trustedNotaryKeys: [NOTARY_PUBLIC_KEY],
  expectedContext: challenge,
});
\end{lstlisting}

\texttt{res} is a standard \texttt{Response}; the proof lives on
\texttt{res.zk}. Sessions serialise (\texttt{toJSON} / \texttt{restoreResponse})
so disclosure can be decided long after the connection closes.

\section{Session modes}

Both modes attest the same object. They differ in who touches the network and
who knows the keys.

\subsection{MPC mode}

In MPC mode $P$ dials $S$ directly and $P$ and $N$ run the TLS client jointly
using two-party computation (DEAP garbled circuits, oblivious transfer, Ferret
VOLE preprocessing).
Session keys are secret-shared: neither party alone knows them, and plaintext
stays with $P$.
The cost is preprocessing of roughly 66~MB, most of it before the request
(\texttt{prepare()} moves it earlier still).
Soundness does not depend on the network path: forging would require breaking
the 2PC or the record authentication.

Browsers cannot open TCP sockets, so MPC mode in a browser needs a
WebSocket-to-TCP relay configured as \texttt{relayUrl}.

\subsection{Proxy mode}

In proxy mode $N$ dials $S$ (port 443) and relays $P$'s TLS bytes, recording
the ciphertext.
$P$ runs a normal TLS client locally, so it knows the keys; afterwards $P$
proves the TLS key schedule in zero knowledge and $N$, for TLS~1.3, decrypts
and checks the handshake flight it relayed (Finished messages, certificate
chain, server name).
Traffic drops to kilobytes instead of tens of megabytes, and browsers need no
relay because $N$ dials.

The trade-off is trust: proxy mode additionally assumes nobody tampers with
the $N \leftrightarrow S$ path, since $P$ knows the server write key and could
otherwise forge records.
The notary therefore dials only public addresses, pins and signs the dialed IP
(verifiers reject a different server name), bounds relayed bytes, and closes
idle connections.
High-value verifiers can set \texttt{rejectProxy} and accept only MPC sessions.

TLS version selection is explicit and never retried: \texttt{auto} means
TLS~1.3 (AES-128-GCM, P-256); TLS~1.2 is available in proxy-compatible
configurations for legacy servers.

\section{What is proven}

\subsection{Key schedule}

The heart of the protocol is proving that the keys used to read the response
are the keys the handshake produced.
In proxy TLS~1.3 sessions $P$ proves the full schedule inside a QuickSilver
zero-knowledge VM: ECDHE secret through HKDF-Extract and Expand-Label to
handshake and application traffic secrets, keys, and IVs.
The default flow uses an ORIGO-minimal schedule that proves only the
compressions touching secret material (about 16 SHA-256 compressions instead
of 32) and reveals one-way intermediate states; a conservative flag restores
the full schedule.
Either way the relation binds every application key and IV to the
server-authenticated handshake transcript.

\subsection{Records and completeness}

TLS records are authenticated encryption with associated data
(AES-128-GCM by default).
$P$ proves, per touched record, the inner content type and the boundary
between content, type byte, and padding, so alerts or tickets cannot be passed
off as HTTP data.
In proxy mode no GCM tag proof is needed ($P$ knows the key; integrity comes
from $N$'s recording plus key binding). In MPC mode tags are checked in 2PC.

A response is marked \texttt{complete} when framing proves it:
\texttt{Content-Length} with matching body length, terminal-chunk framing, or
a recorded \texttt{close\_notify} to the body end.
Strict verifiers require completeness; anything else is a proven prefix, and
applications must treat it as such.

\subsection{Commitments}

Commitments bind the transcript without revealing it.
By default each header and each JSON node gets a hash commitment computed
inside the ZK VM, and the notary signs them.
An important optimisation, \texttt{zkConfig.reveal}, declares the intended
disclosure up front: $P$ then commits to one BLAKE3 commitment per disclosed
unit (header, JSON field, target, body, skeleton) and commits nothing for
hidden values.
This removes about 54\% of the ZK work (1.62M of 3.0M AND gates per session
in the measured fixture): proxy proving falls from about 135~ms to 59~ms and
presentations shrink from 22.8~KB to 6.7~KB.
The restriction is enforced: later presentations may disclose that spec or
less (by whole units), never more. The opt-in Binius backend can still prove
new predicates later.

\subsection{Predicates and JSON disclosure}

Two backends prove facts about hidden values.
\textbf{QuickSilver} (default) proves numeric predicates (\texttt{gte},
\texttt{gt} over unsigned integers up to 19 digits) to $N$ during the session;
$N$ signs the results, so offline verification is a signature check.
\textbf{Binius64} (opt-in) adds per-leaf BLAKE3 commitments so $P$ can prove
new predicates later, offline, without $N$ --- at the cost of 734--760~KiB
proofs and hundreds of milliseconds.

JSON fields are selected by dotted path (e.g.\ \texttt{users.0.username}),
including array elements.
With QuickSilver the notary attests the shape of every JSON value, so
\texttt{verify()} returns each disclosed value with a proven path.
Disclosing a field also discloses the response's keys and punctuation (the
skeleton), but not other values.
A byte-only mode skips the skeleton and returns raw spans without path
authentication; sessions are capped at 1,024 leaves, and ambiguous paths
(empty keys, keys containing periods, duplicate decoded keys) are rejected.

\section{Verification}

Offline verification is deliberately boring:

\begin{enumerate}[leftmargin=2em]
\item Decode strictly; reject unknown critical fields.
\item Check $N$'s signature against a pinned key within its validity window;
fail closed on expired revocation data.
\item Check policy: server name, mode, method, target, status, freshness
(\texttt{maxAgeSecs} on the verifier's clock), \texttt{context},
\texttt{owner}, required predicates, completeness, and path authentication.
\item Check each commitment opening and each predicate proof.
\item Return the redacted transcript (hidden bytes shown as \texttt{X}),
the authenticated byte ranges, the proven JSON fields, and the predicate list.
\end{enumerate}

Applications must decide from the authenticated ranges
(\texttt{recvAuthed} / \texttt{sentAuthed}), never from the display strings,
since a hidden byte and a literal \texttt{X} look identical on screen.

\section{Security model}

\textbf{Against a malicious prover,} data integrity and server authenticity
hold given an honest notary (plus, in proxy mode, an honest $N \leftrightarrow
S$ path): forging a response, predicate, or server identity requires breaking
the key-schedule proof, the record proofs, the certificate validation, or the
signature scheme.
\textbf{Against a malicious notary,} privacy holds: $N$ sees ciphertext,
lengths, one-way handshake states, IVs, and zero-knowledge transcripts, but
never plaintext or application keys.
\textbf{Against a malicious verifier,} presentations reveal only the selected
bytes, block positions, total length, and one bit per predicate; replay is
stopped by the signed \texttt{context} challenge.
$P$ and $N$ together can always forge; no single-notary protocol avoids this
without additional trust, so deployments that need it run the notary
themselves (verifier-as-notary) and publish keys and attestations in a
transparency log.

Known limits are stated openly: the TLS~1.3 and proxy constructions have not
had an independent malicious-security review; chunked and EOF-delimited bodies
satisfy only prefix authentication under the strict completeness policy;
JSON disclosure leaks structure and lengths; Binius lacks a reviewed soundness
argument for this use; and several dependencies are flagged for replacement.
The implementation note \texttt{ZKFETCH\_PATCHES.md} lists every change to the
vendored TLSNotary and mpz stacks.

\section{Performance}

All optimisations keep every wire byte identical, so old and new provers and
notaries interoperate.
Representative measurements (local fixture on an M2 Max; live Duolingo
from India to a Mumbai \texttt{t4g.small} notary; browser single-threaded
WebAssembly):

\begin{center}
\begin{tabular}{ll}
\toprule
Session & Time \\
\midrule
Local proxy, native (fixture) & $\sim$185~ms total, $\sim$112~ms proving \\
Local proxy with declared reveal & $\sim$130~ms total, $\sim$59~ms proving \\
Live proxy, native (Duolingo) & $\sim$2.1~s (was 15--17~s via Cloudflare) \\
Live proxy, browser, unprepared & $\sim$15~s (connect 0.9, setup 6.9, TLS 1.5, proofs 5.1, attest 0.6) \\
Live proxy, browser, prepared & 2.3--3.0~s after click \\
\bottomrule
\end{tabular}
\end{center}

Latency is dominated by about 21 sequential prover--notary round trips plus
VOLE preprocessing.
Persistent Ferret bootstrap pools (single-use, transcript-bound leases kept in
memory across calls) cut setup from 513~ms to 138~ms and end-to-end
notarization from 963~ms to 585~ms at 33~ms RTT, with 40\% less traffic.
A three-exchange session flow and the ORIGO-minimal schedule further cut warm
medians substantially.
ARM64 PMULL multiplication, an eight-row LPN encoder, lazy field reduction,
WebAssembly SIMD, and eight-thread proving (3.3~s $\rightarrow$ 1.1~s prover
CPU) compound the gains.
QuickSilver presentations are 14--24~KiB with sub-millisecond verification;
Binius presentations are two orders of magnitude larger.

\section{Related work and outlook}

zkFetch descends from DECO and TLSNotary: MPC-TLS with garble-then-prove for
the handshake, now extended with a relayed proxy mode, TLS~1.3 proofs, and
JSON-shape circuits.
The planned v2 architecture keeps proxy mode with ZK-proven key derivation
(the construction that recent analyses by Luo et al.\ and ORIGO show to be
sound), commits to keys rather than plaintext, minimises the key-schedule
circuit, collapses the message flow to three round trips, replaces record-layer
MPC with a key-schedule-only split mode, and swaps Binius64 for
VOLE-in-the-Head public proofs --- predicting roughly a further halving of
session time and twenty-fold smaller offline proofs.
Post-quantum key exchange, anchored JSON disclosure without skeleton leakage,
and an optional on-chain signed-claim verifier follow the same pattern:
prove less, in zero knowledge, and sign exactly what was checked.

\section{Conclusion}

zkFetch shows that verifiable web data does not need new servers, new hardware,
or all-or-nothing disclosure.
A notary that dials or co-signs, a prover that proves keys rather than trust,
and a verifier that checks signatures and small proofs are enough to turn any
HTTPS response into evidence: reveal a name, prove a streak, hide the token.
The protocol is implemented, measured, and deployed on commodity infrastructure;
its patches, limits, and next steps are public.

\end{document}
